\chapter{Bibliography}

\section*{Primary Sources}

\begin{description}[style=nextline, font=\normalfont]

\item[\textit{Nettippakaraṇa}.] Chaṭṭha Saṅgāyana Tipiṭaka Version 4.0. Vipassana Research Institute, Igatpuri. Digital Pali text.

\item[\textit{Nettippakaraṇa}.] Pali Text Society edition. Ed. E. Hardy. London: PTS, 1902.

\item[\textit{The Guide} (\textit{Netti-pakaraṇa}).] Trans. Bhikkhu Ñāṇamoli. London: Pali Text Society, 1962. Reprinted 1977, 2015.

\item[\textit{Peṭakopadesa}.] Pali Text Society edition. Ed. A. Barua. London: PTS, 1949. The ``sister text'' of the Netti, covering much of the same hermeneutical system.

\item[\textit{Peṭakopadesa: Pitaka-Disclosure}.] Trans. Bhikkhu Ñāṇamoli. London: Pali Text Society, 1964.

\end{description}

\section*{Canonical Collections Cited}

\begin{description}[style=nextline, font=\normalfont]

\item[\textit{Dīgha Nikāya} (DN).] ``Collection of Long Discourses.'' Trans. Maurice Walshe as \textit{The Long Discourses of the Buddha}. Boston: Wisdom Publications, 1995.

\item[\textit{Majjhima Nikāya} (MN).] ``Collection of Middle-Length Discourses.'' Trans. Bhikkhu Ñāṇamoli and Bhikkhu Bodhi as \textit{The Middle Length Discourses of the Buddha}. Boston: Wisdom Publications, 1995.

\item[\textit{Saṃyutta Nikāya} (SN).] ``Collection of Connected Discourses.'' Trans. Bhikkhu Bodhi as \textit{The Connected Discourses of the Buddha}. Boston: Wisdom Publications, 2000.

\item[\textit{Aṅguttara Nikāya} (AN).] ``Collection of Numerical Discourses.'' Trans. Bhikkhu Bodhi as \textit{The Numerical Discourses of the Buddha}. Boston: Wisdom Publications, 2012.

\item[\textit{Khuddaka Nikāya}.] ``Collection of Minor Texts.'' Including the Dhammapada, Udāna, Itivuttaka, Sutta Nipāta, Theragāthā, Therīgāthā, and Jātaka, among others. Individual translations cited by translator.

\item[\textit{Dhammapada} (Dhp).] Trans. Gil Fronsdal as \textit{The Dhammapada}. Boston: Shambhala, 2005.

\item[\textit{Udāna} (Ud).] Trans. John D. Ireland as \textit{The Udāna: Inspired Utterances of the Buddha}. Kandy: Buddhist Publication Society, 1997.

\item[\textit{Sutta Nipāta} (Sn).] Trans. Bhikkhu Bodhi as \textit{The Suttanipāta}. Boston: Wisdom Publications, 2017.

\end{description}

\section*{Secondary Sources}

\begin{description}[style=nextline, font=\normalfont]

\item[Bodhi, Bhikkhu.] \textit{The Noble Eightfold Path: Way to the End of Suffering}. Kandy: Buddhist Publication Society, 1994.

\item[Bodhi, Bhikkhu.] ``The Netti-pakaraṇa: An Ancient Hermeneutics of the Pali Canon.'' In \textit{Encyclopaedia of Buddhism}, ed. G.P. Malalasekera. Colombo: Government of Sri Lanka, 2003.

\item[Bond, George D.] ``The Nature and Meaning of the Nettippakaraṇa.'' In \textit{Buddhist Studies in Honour of Walpola Rahula}, ed. S. Balasooriya et al. London: Gordon Fraser, 1980.

\item[Cousins, Lance S.] ``Pali Oral Literature.'' In \textit{Buddhist Studies: Ancient and Modern}, ed. P. Denwood and A. Piatigorsky. London: Curzon Press, 1983.

\item[Gethin, Rupert.] \textit{The Foundations of Buddhism}. Oxford: Oxford University Press, 1998.

\item[Ñāṇamoli, Bhikkhu.] \textit{The Path of Purification} (\textit{Visuddhimagga} of Bhadantācariya Buddhaghosa). Kandy: Buddhist Publication Society, 1991.

\item[Warder, A.K.] \textit{Indian Buddhism}. Delhi: Motilal Banarsidass, 1970. Revised edition 2000.

\end{description}

\section*{Digital Resources}

\begin{description}[style=nextline, font=\normalfont]

\item[SuttaCentral (suttacentral.net).] Parallel editions and translations of early Buddhist texts, including the Nettippakaraṇa. Cited for cross-references to Nikāya parallels.

\item[Digital Pali Reader.] Software for reading and searching the Chaṭṭha Saṅgāyana Tipiṭaka.

\item[Tipitaka XML (tipitaka.org).] XML-encoded Pali texts used as the primary digital source for this translation.

\end{description}
